\section{省略证明}

\subsection{定理~\ref{thm:pivot} 的证明}

\begin{proof}
固定状态 $s$ 和 $\beta>0$。为简便起见，记
\begin{align*}
\pi = \pi_{s,\beta}, \quad
q = \mathbb E_{a\sim\pi}[r(s,a)], \quad
\sigma^2 = \mathrm{Var}_{a\sim\pi}(r(s,a)).
\end{align*}

首先计算 $J_s$ 在 $\pi$ 处的自然梯度。令 $v\in T_\pi$，考虑指数扰动
\begin{align*}
\pi_t(a)= &~ \frac{\pi(a)e^{t v(a)}}{\sum_{b\in\mathcal{A}(s) }\pi(b)e^{t v(b)}}.
\end{align*}
对任意 $v\in T_\pi$，有
\begin{align*}
\left.\frac{d}{dt}\right|_{t=0}\log \pi_t(a)
= v(a)-\sum_{b\in\mathcal{A}(s) }\pi(b)v(b)
= v(a).
\end{align*}
因此
\begin{align*}
DJ_s(\pi)[v]
= &~ \left.\frac{d}{dt}\right|_{t=0}\mathbb E_{a\sim\pi_t}[r(s,a)] \\
= &~ \left.\frac{d}{dt}\right|_{t=0}\sum_{a\in\mathcal{A}(s) }\pi_t(a)r(s,a) \\
= &~ \sum_{a\in\mathcal{A}(s) }\pi(a)r(s,a)v(a) \\
= &~ \sum_{a\in\mathcal{A}(s) }\pi(a)\bigl(r(s,a)-q\bigr)v(a) \\
= &~ \langle r(s,\cdot)-q,\,v\rangle_{F,\pi}.
\end{align*}
由于这一结论对所有 $v\in T_\pi$ 成立，其 Riesz 表示元为
\begin{align*}
\nabla^{\mathrm{nat}}J_s(\pi)= &~ r(s,\cdot)-q.
\end{align*}
于是有
\begin{align}\label{eq:norm-natural-J}
\bigl\|\nabla^{\mathrm{nat}}J_s(\pi)\bigr\|_{F,\pi}^2
= &~ \mathbb E_{a\sim\pi}\bigl[(r(s,a)-q)^2\bigr] \notag\\
= &~ \mathrm{Var}_{a\sim\pi}(r(s,a)) \notag\\
= &~ \sigma^2.
\end{align}

接着计算指数倾斜路径对 $\beta$ 的导数。定义
\begin{align*}
Z_s(\beta)= &~ \sum_{b\in\mathcal{A}(s) }\pi_0(b\mid s)e^{r(s,b)/\beta}.
\end{align*}
则
\begin{align*}
\log \pi_{s,\beta}(a)
= &~ \log \pi_0(a\mid s) + \frac{r(s,a)}{\beta} - \log Z_s(\beta).
\end{align*}
对 $\beta$ 求导得
\begin{align*}
\partial_\beta \log \pi_{s,\beta}(a)
= &~ -\frac{r(s,a)}{\beta^2} - \frac{Z_s'(\beta)}{Z_s(\beta)}.
\end{align*}
而
\begin{align*}
Z_s'(\beta)
= &~ \sum_{b\in\mathcal{A}(s) }\pi_0(b\mid s)e^{r(s,b)/\beta}\left(-\frac{r(s,b)}{\beta^2}\right) \\
= &~ -\frac{1}{\beta^2}\sum_{b\in\mathcal{A}(s) }\pi_0(b\mid s)e^{r(s,b)/\beta}r(s,b) \\
= &~ -\frac{Z_s(\beta)}{\beta^2}\,\mathbb E_{b\sim\pi_{s,\beta}}[r(s,b)] \\
= &~ -\frac{q}{\beta^2}Z_s(\beta).
\end{align*}
因此
\begin{align*}
\partial_\beta \log \pi_{s,\beta}(a)
= &~ -\frac{1}{\beta^2}\bigl(r(s,a)-q\bigr) \\
= &~ -\frac{1}{\beta^2}\nabla^{\mathrm{nat}}J_s(\pi_{s,\beta})(a).
\end{align*}

若 $\sigma=0$，则对 $\pi$-几乎处处的 $a$ 都有 $r(s,a)=q$，于是
\begin{align*}
\nabla^{\mathrm{nat}}J_s(\pi)= &~ 0, ~~\text{并且}~~ \frac{1}{\beta^2}\bigl\|\nabla^{\mathrm{nat}}J_s(\pi)\bigr\|_{F,\pi}= 0.
\end{align*}

现在假设 $\sigma>0$。则
\begin{align*}
\gamma_{s,\beta}
= &~ -\mathbb E_{a\sim\pi}
\left[
\frac{r(s,a)-q}{\sigma}\,\partial_\beta \log \pi(a)
\right] \\
= &~ -\mathbb E_{a\sim\pi}
\left[
\frac{r(s,a)-q}{\sigma}\left(-\frac{1}{\beta^2}(r(s,a)-q)\right)
\right] \\
= &~ \frac{1}{\beta^2\sigma}\mathbb E_{a\sim\pi}\bigl[(r(s,a)-q)^2\bigr] \\
= &~ \frac{\sigma}{\beta^2}.
\end{align*}
结合式~\eqref{eq:norm-natural-J}，得到
\begin{align*}
\gamma_{s,\beta}
= &~ \frac{1}{\beta^2}\bigl\|\nabla^{\mathrm{nat}}J_s(\pi_{s,\beta})\bigr\|_{F,\pi_{s,\beta}} = \frac{\sigma}{\beta^2}.
\end{align*}
定理得证。
\end{proof}

\subsection{定理~\ref{thm:kl} 的证明}\label{sec:proof-kl-proj}

\begin{proof}
由于
\begin{align*}
\mathcal L_{\mathrm{func},\beta}(\pi)= &~ \mathbb E_{s\sim d}[\ell_s(\pi)]
\end{align*}
在不同状态之间没有耦合，因此只需对每个固定状态 $s$ 求解逐点极小化问题。总体目标的结论随即在 $d$-几乎处处意义下成立。

固定状态 $s$，记 $M= \mathcal M(s), ~\rho= \pi_0(M\mid s), ~\pi_0(a)= \pi_0(a\mid s)$。

我们先处理边界情形。若 $\rho=0$，则由于对每个 $a\in \mathcal{A}(s)$ 都有 $\pi_0(a)>0$，可知 $M=\varnothing$。于是 $r_{\mathrm{func}}(s,a)\equiv 0$，从而 $\ell_s(\pi)= \beta\, \mathrm{KL}(\pi\|\pi_0)$，
其唯一极小值在 $\pi=\pi_0$ 处取得。若 $\rho=1$，则 $M=\mathcal{A}(s)$，因此 $r_{\mathrm{func}}(s,a)\equiv 1$，并有 $\ell_s(\pi)= -1 + \beta\, \mathrm{KL}(\pi\|\pi_0)$，
其唯一极小值同样在 $\pi=\pi_0$ 处取得。

接下来考虑非退化情形 $0<\rho<1$。令 $p$ 为 $\mathcal{A}(s)$ 上任意一个分布，并定义
\begin{align*}
q= &~ p(M)=\sum_{a\in M} p(a).
\end{align*}
对 $q\in(0,1)$，定义条件分布
\begin{align*}
p_+(a)= \frac{p(a)}{q}, \quad a\in M, \qquad
p_-(a)= \frac{p(a)}{1-q}, \quad a\notin M,
\end{align*}
类似地，
\begin{align*}
\pi_{0,+}(a)= \frac{\pi_0(a)}{\rho}, \quad a\in M, \qquad
\pi_{0,-}(a)= \frac{\pi_0(a)}{1-\rho}, \quad a\notin M.
\end{align*}

利用 KL 散度的链式法则，可得
\begin{align*}
\mathrm{KL}(p\|\pi_0)
= &~ \sum_{a\in M} p(a)\log\frac{p(a)}{\pi_0(a)}
   + \sum_{a\notin M} p(a)\log\frac{p(a)}{\pi_0(a)} \\
= &~ q\,\mathrm{KL}(p_+\|\pi_{0,+})
   + (1-q)\,\mathrm{KL}(p_-\|\pi_{0,-}) + q\log\frac{q}{\rho}
   + (1-q)\log\frac{1-q}{1-\rho}.
\end{align*}

又因为 $\mathbb E_{a\sim p}[r_{\mathrm{func}}(s,a)]=q$，逐点目标可写为
\begin{align*}
\ell_s(p)
= -q
+ \beta\Bigg[
q\,\mathrm{KL}(p_+\|\pi_{0,+})
+ (1-q)\,\mathrm{KL}(p_-\|\pi_{0,-})
+ q\log\frac{q}{\rho}
+ (1-q)\log\frac{1-q}{1-\rho}
\Bigg].
\end{align*}

在固定 $q$ 时，最后两项是常数，而前两项 KL 项非负，并且当且仅当
\begin{align*}
p_+= \pi_{0,+}, \quad
p_-= \pi_{0,-}
\end{align*}
时取零。因此，在固定 $q$ 时，唯一极小化解为
\begin{align*}
\bar p_q(a)= &~
\begin{cases}
\displaystyle \frac{q}{\rho}\,\pi_0(a), & a\in M,\\[1.1ex]
\displaystyle \frac{1-q}{1-\rho}\,\pi_0(a), & a\notin M.
\end{cases}
\end{align*}

把 $p=\bar p_q$ 代回去，条件 KL 项被消去，得到如下标量问题
\begin{align*}
\phi_s(q)= &~ -q
+ \beta\left[
q\log\frac{q}{\rho}
+ (1-q)\log\frac{1-q}{1-\rho}
\right],
\qquad q\in[0,1],
\end{align*}
其中采用惯例 $0\log 0 = 0$。

当 $q\in(0,1)$ 时，
\begin{align*}
\phi_s'(q)= &~ -1 + \beta\log\frac{q(1-\rho)}{(1-q)\rho}, \\
\phi_s''(q)= &~ \beta\left(\frac{1}{q}+\frac{1}{1-q}\right) > 0.
\end{align*}
因此 $\phi_s$ 在 $(0,1)$ 上严格凸，且其中至多只有一个驻点。解 $\phi_s'(q)=0$ 得
$q= \frac{\rho e^{1/\beta}}{(1-\rho)+\rho e^{1/\beta}}$。
于是，$\ell_s$ 的唯一极小值为 $\pi_\beta^\star(\cdot\mid s)=\bar p_{q_\beta(s)}$，即
\begin{align}\label{eq:pi-beta-star}
\pi_\beta^\star(a\mid s)= &~
\begin{cases}
\displaystyle \frac{q_\beta(s)}{\rho(s)}\,\pi_0(a\mid s), & a\in \mathcal M(s),\\[1.1ex]
\displaystyle \frac{1-q_\beta(s)}{1-\rho(s)}\,\pi_0(a\mid s), & a\notin \mathcal M(s).
\end{cases}
\end{align}
在所有满足可接受动作总质量为 $q_\beta(s)$ 的分布中，上述分块缩放后的分布正是唯一的 KL 极小化解。

由式~\eqref{eq:pi-beta-star}，排序保持性质立刻成立。若 $a,b\in \mathcal M(s)$，则
\begin{align*}
\frac{\pi_\beta^\star(a\mid s)}{\pi_\beta^\star(b\mid s)}
= &~ \frac{\pi_0(a\mid s)}{\pi_0(b\mid s)}.
\end{align*}
同理，若 $a,b\notin \mathcal M(s)$，则
\begin{align*}
\frac{\pi_\beta^\star(a\mid s)}{\pi_\beta^\star(b\mid s)}
= &~ \frac{\pi_0(a\mid s)}{\pi_0(b\mid s)}.
\end{align*}
证明完毕。
\end{proof}
